\documentclass[11pt,twoside]{book}
\usepackage[T1]{fontenc}
\usepackage[utf8]{inputenc}
\usepackage{lmodern}
\usepackage[margin=1in,bottom=1.15in]{geometry}
\usepackage{amsmath,amssymb,mathtools}
\usepackage{fancyhdr}
\usepackage{titlesec}
\usepackage{xcolor}
\usepackage{wrapfig}
\usepackage{tikz}
\usepackage[hidelinks]{hyperref}
\definecolor{physicsred}{HTML}{8C2F39}
\pagestyle{fancy}
\fancyhf{}
\fancyhead[LE]{\thepage\quad\small\nouppercase{\leftmark}}
\fancyhead[RO]{\small\nouppercase{\rightmark}\quad\thepage}
\renewcommand{\headrulewidth}{0pt}
\titleformat{\chapter}[display]{\normalfont\bfseries\Huge\color{physicsred}}
  {\filleft\MakeUppercase{\chaptertitlename}~\thechapter}{12pt}{\filright}
\begin{document}
\begin{titlepage}
\centering
\vspace*{0.8in}
{\fontsize{38}{44}\selectfont\bfseries Spacetime\\and Motion\par}
\vspace{18pt}\textcolor{physicsred}{\rule{0.65\textwidth}{2pt}}\par
\vspace{20pt}{\Large Notes in introductory physics\par}
\vfill
{\Large Author Name\par}\vspace{6pt}{\large Institution or laboratory\par}
\vspace{0.5in}{\small First edition \textbullet{} \the\year\par}
\end{titlepage}
\frontmatter
\chapter*{How to use this book}
Equations, short explanations, and margin observations are arranged to make
the mathematics visible. Replace the sample notes with the conventions and
examples of your own course or monograph.
\tableofcontents
\mainmatter
\chapter{Frames of reference}
\section{Describing an event}
An event is specified by a position and a time. Different observers may use
different coordinates, but they can still agree on the event itself.
\begin{equation}
  \mathbf r(t)=x(t)\,\mathbf e_x+y(t)\,\mathbf e_y+z(t)\,\mathbf e_z.
\end{equation}
\section{A physical estimate}
Before calculating, estimate the scale of the answer. Dimensional analysis is
often the quickest way to find an incorrect factor or an impossible unit.
\begin{equation}
  [v]=\frac{\mathrm{length}}{\mathrm{time}},\qquad
  [a]=\frac{\mathrm{length}}{\mathrm{time}^2}.
\end{equation}
\chapter{Dynamics}
\section{A compact derivation}
For constant acceleration, integrating $a=\mathrm{d}v/\mathrm{d}t$ twice gives
\begin{align}
v(t)&=v_0+at,\\
x(t)&=x_0+v_0t+\tfrac12at^2.
\end{align}
\begin{quote}
\textbf{Observation.} A derivation is easier to check when each line preserves
the dimensions of the physical quantity it defines.
\end{quote}
\chapter{Problems}
\begin{enumerate}
  \item Choose a coordinate system for a projectile and state its origin.
  \item Derive the time at which the projectile reaches its highest point.
  \item Which measurements would most strongly affect the final estimate?
\end{enumerate}
\backmatter
\chapter*{Notation}
Use this page for symbols, constants, and conventions shared across chapters.
\end{document}
